\textbf{Distance calculation}

The information from the table can be used to derive individual magnitudes of
both components according to equations:
\begin{align*}
  m_1 &= m_{1+2} + 2.5\log_{10}(1 + L_2/L_1) \\
  m_2 &= m_{1+2} + 2.5\log_{10}(1 + L_1/L_2).
\end{align*}
We will use the third system (OGLE LMC-ECL-18365) as an example to
demonstrate the calculations in detail. The values for this system are as
follows:

Magnitudes: $m_{V,1} = 16.47$\,mag, $m_{K,1} = 14.20$\,mag,
$m_{V,2} = 18.19$\,mag, $m_{K,2} = 16.01$\,mag.

Colours: $(m_V - m_K)_1 = 2.27$\,mag, $(m_V - m_K)_2 = 2.18$\,mag.

The second step is to determine the surface brightness $S_V$ for both
components using the figure. A least square fit to the data in the figure
results in a linear function:
\begin{align*}
  S_V &= 1.346\,\big((m_V - m_K) - 2.407\big) + 5.869\ [\mathrm{mag}],
    \text{ or} \\
  S_V &= 1.346\,(m_V - m_K) + 2.629\ [\mathrm{mag}].
\end{align*}
Uncertainties of the coefficients are $1.346 \pm 0.017$\,mag and
$2.63 \pm 0.04$\,mag.

For this system, this gives $S_{V,1} = 5.69$\,mag and $S_{V,2} = 5.57$\,mag.

However, participants will have two other ways of determining $S_V$.

1) The first way is a graphical way by using a ruler and a pencil in order to
draw the `best-fit' line on the figure. Then $S_V$ follows from an
intersection of the $x = (m_V - m_K)$ vertical line and the `best-fit' line.

2) The second way is to determine the coefficients of the best-fit line
$y = ax + b$ by using coordinates of two points on the figure. The points
should be far from each other.

For example, the coordinates of the second and the penultimate points are:
$(x_1, y_1) = (2.07, 5.41)$ and $(x_2, y_2) = (2.71, 6.28)$. This results in:
\begin{align*}
  a &= \frac{y_2 - y_1}{x_2 - x_1} = 1.36 \\
  b &= \frac{y_1 x_2 - y_2 x_1}{x_2 - x_1} = 2.60
\end{align*}
Using these coefficients we have:
$S_{V,1} = y = 2.27\cdot a + b = 5.69$\,mag and
$S_{V,2} = y = 2.18\cdot a + b = 5.56$\,mag, so within 0.01\,mag of the
`precise' results.

The third step is to calculate angular diameters of components by using the
equation presented in the problem. By modifying the equation defining $S_V$
we obtain:
\[ \theta = 10^{0.2(S_V - m_V)} \]
Subsequently we get: $\theta_1 = 10^{0.2(5.69-16.47)} = 0.00698$\,mas and
$\theta_2 = 10^{0.2(5.56-18.19)} = 0.00298$\,mas.

The fourth step is to calculate the distance to each target. As components
form a physical binary, their distances should be very similar. This is an
independent check of the method and calculations. As the angles under which
we see stellar discs are very small ($\sin\theta \approx \theta$) we can
safely use a linear relation between the angular and physical diameters of an
object. We therefore calculate the distance $D$ as $D = kR/\theta$, where
$\theta$ is expressed in mas, $R$ in solar radii and $k$ is a conversion
factor. The conversion factor results from a fraction
$(2R_\odot/1\,\mathrm{kpc})/(1\,\mathrm{mas}/1\,\mathrm{rad})
= (2R_\odot/1\,\mathrm{AU}) = 9.30\cdot 10^{-3}$.

The distances for the third system are
$D_1 = 9.30\cdot 10^{-3}\cdot 37.30/0.00698 = 49.70$\,kpc and
$D_2 = 9.30\cdot 10^{-3}\cdot 15.94/0.00298 = 49.75$\,kpc.

We then repeat the calculation following the same scheme for the first and
second systems, obtaining for the first system $D_1 = 49.33$\,kpc and
$D_2 = 49.07$\,kpc, and for the second system $D_1 = 49.56$\,kpc and
$D_2 = 49.30$\,kpc. The unweighted mean of all distances is 49.45\,kpc.

\medskip
\textbf{Uncertainties}

\emph{`Statistical' part.}

The standard deviation of the sample is $s = 0.24$\,kpc. The standard error
of the mean is $s/\sqrt{6} = 0.09$\,kpc.

OR:

The mean distances to the three eclipsing binaries are: 49.73\,kpc,
49.20\,kpc and 49.443\,kpc. The standard deviation is $s = 0.22$\,kpc, and
the standard error of the mean is $s/\sqrt{3} = 0.13$\,kpc.

\emph{`Systematic' part.}

All distances are inversely proportional to angular diameters derived from
the $S_V$\,--\,$(m_V - m_K)$ relation. Thus their accuracy is limited by the
precision of the relation. That gives the `irreducible' part of the error:
$0.008\cdot 49.45 = 0.40$\,kpc.

Finally: the distance is $49.45 \pm 0.09 \pm 0.40$\,kpc; the uncertainty is
completely dominated by the precision of the $S_V$\,--\,$(m_V - m_K)$
relation.
